\originalchapter{密度扰动为何增长，以及何时停止增长}
本章对应8.902第12讲末、第13讲和第14讲的增长部分。星系不是均匀Friedmann解本身的产物，而来自小的密度起伏。我们先在适用的尺度上从流体方程推导增长，再说明辐射、本动速度和规范的限制。原笔记把一个启发式辐射方程延伸到视界外，本卷保留其物理主题，改用有明确适用条件的论述。

\originalsection{在膨胀背景中定义扰动}
令$\rho(\mathbf x,t)=\bar\rho(t)[1+\delta(\mathbf x,t)]$，$\mathbf r=a\mathbf x$，物理速度$\mathbf v=H\mathbf r+\mathbf u$。$\delta$是相对于同一时刻背景的密度对比，不等于绝对密度差；$\mathbf u$是本动速度。对非相对论、亚Hubble尺度、弱引力势的流体，原物理坐标方程变换到共动坐标并减去背景后给出
\begin{align}
 \dot\delta+\frac1a\boldsymbol\nabla\!\cdot[(1+\delta)\mathbf u]&=0,\label{eq:astro-pert-cont}\\
 \dot{\mathbf u}+H\mathbf u+\frac1a(\mathbf u\cdot\boldsymbol\nabla)\mathbf u
 &=-\frac1a\boldsymbol\nabla\phi-\frac{\boldsymbol\nabla p}{a\rho},\label{eq:astro-pert-euler}\\
 \nabla^2\phi&=4\pi Ga^2\bar\rho\delta.\label{eq:astro-pert-poisson}
\end{align}
梯度作用于$\mathbf x$。$H\mathbf u$可理解为自由粒子的本动物理动量随$a^{-1}$衰减。引力加速度是$-\boldsymbol\nabla\phi/a$；原笔记的正号与其$\nabla^2\phi=4\pi G\rho$不一致。

暗物质原则上由无碰撞Boltzmann方程描述；上述流体式是其矩方程，在速度分布足够冷、尚未壳层交叉时可近似为单流尘埃。发生多流时，平均速度仍存在，但不能忽略速度弥散应力，也不能把不同流简单当成一个无压强速度场。

\originalsection{线性方程的两个Hubble阻尼因子}
保留$|\delta|\ll1$和$\mathbf u$的一阶项，取绝热压强扰动$\delta p=c_s^2\bar\rho\delta$，有$\dot\delta=-\nabla\cdot\mathbf u/a$。求时间导数时，一个$H$来自$1/a$，另一个来自Euler方程中的$H\mathbf u$，代入Poisson式得
\begin{equation}\label{eq:astro-growth}
 \ddot\delta+2H\dot\delta-4\pi G\bar\rho\delta
 =\frac{c_s^2}{a^2}\nabla^2\delta.
\end{equation}
这一步不能漏掉共动梯度前的$a^{-1}$，否则会错误地得到只有一次$H$阻尼。对
$\delta(\mathbf x,t)=\int\delta_{\mathbf k}(t)\ee^{\ii\mathbf k\cdot\mathbf x}\dd^3k/(2\pi)^3$，方程成为
\[
 \ddot\delta_{\mathbf k}+2H\dot\delta_{\mathbf k}
 +\left(\frac{c_s^2k^2}{a^2}-4\pi G\bar\rho\right)\delta_{\mathbf k}=0.
\]
物理波数为$k/a$。瞬时Jeans分界$k_{J,\rm phys}=\sqrt{4\pi G\bar\rho}/c_s$，物理波长$\lambda_J=c_s\sqrt{\pi/(G\bar\rho)}$。膨胀使系数随时间变化，不能把增长时间始终当成静态指数$\ee^{t\sqrt{4\pi G\rho}}$。

\originalsection{物质主导时期的两个独立解}
Einstein--de Sitter背景有$H=2/(3t)$、$\bar\rho=1/(6\pi Gt^2)$。忽略压强，试$\delta\propto t^n$，式\eqref{eq:astro-growth}给
\[
 n(n-1)+\frac43n-\frac23=0,
 \qquad n=\frac23,-1.
\]
所以增长模$D_+\propto a$，衰减模$D_-\propto t^{-1}=a^{-3/2}$。线性叠加的系数由初始密度与速度共同决定；只指定$\delta(t_i)$不足以唯一决定增长。常用增长因子$D(a)$约定$D(1)=1$，不是把每个模的今天振幅都设为1。

暗能量或曲率开始影响背景时，冷物质自引力项仍含$\bar\rho_m$，而$H$包含所有背景组分。用$a$作自变量，记$D_a=\dd D/\dd a$，得
\[
 D_{aa}+\left(\frac3a+\frac{E_a}{E}\right)D_a
 -\frac{3\Omega_{m,0}}{2a^5E^2}D=0.
\]
在忽略辐射、包含物质、曲率和恒定$\Lambda$的模型中，$E(a)$本身是一个独立解。降阶法使另一个解为$E(a)\int\dd a/(a^3E^3)$，取早期$D_+\simeq a$的归一化有
\[
 \widetilde D_+(a)=\frac{5\Omega_{m,0}}2E(a)\int_0^a\frac{\dd a'}{a'^3E^3(a')},
 \qquad D(a)=\frac{\widetilde D_+(a)}{\widetilde D_+(1)}.
\]
正$\Lambda$占优势时，密度下降而$H$趋向常数，增长逐渐冻结。这一积分在有不可忽略辐射或一般时变暗能量时不直接适用，应解原方程。

\originalsection{光子与重子为何会振荡}
复合前，Thomson散射将电子与光子紧耦合，Coulomb作用又将电子与离子耦合。对绝热扰动$\delta\rho_b/\rho_b=(3/4)\delta\rho_\gamma/\rho_\gamma$，定义$R_b=3\rho_b/(4\rho_\gamma)$，总压强由辐射给出$\delta p=c^2\delta\rho_\gamma/3$，于是有效声速
\[
 c_{s,\gamma b}^2=\frac{c^2}{3(1+R_b)}.
\]
只有重子惯性可以忽略时才是$c/\sqrt3$。在因果联系内，光子压强与引力竞争，使扰动进行声学振荡；重子不再能像冷尘埃那样自由增长。复合解除光子压强支持后，气体声速下降到公里每秒量级，能进入已有暗物质势阱。

若把Jeans质量定义为\emph{直径为}$\lambda_J$的均匀球质量，则$M_J=(\pi/6)\rho\lambda_J^3$。原课以半径$\lambda_J$定义，故相差8倍；课程的$10^{16}\Ms$和$10^5\Ms$应理解为相应物理阶段的尺度估计，不能在未统一球半径和温度的情况下比较系数。复合后气体在一段时间内仍通过残余电子的Compton散射与CMB保持热耦合，后来绝热冷却趋于$T_b\propto a^{-2}$，并非一直$T_b=T_\gamma\propto a^{-1}$。

在紧耦合不是无限强时，光子走随机游走，扩散物理距离约为$\sqrt{\lambda_{\rm mfp}ct}$；越近脱耦，平均自由程越大。光子携带热量和动量离开密度峰，阻尼小尺度声学扰动，这称Silk阻尼。更精确的共动阻尼积分要同时计入剪切黏性、热传导和随时间变化的电离率；课程的Silk质量约$10^{12}\Ms$只是数量级，不表示所有更小质量对象永远不能后来形成。

\originalsection{冷暗物质在辐射时期的增长为何受抑制}
冷暗物质不受光子压强直接支持，但其密度在辐射主导背景中较小。对亚Hubble冷物质扰动，若在声学振荡平均后忽略辐射密度扰动的引力源，仍用式\eqref{eq:astro-growth}，只是$H^2\propto\bar\rho_r+\bar\rho_m$而自引力只含$\bar\rho_m$。令$y=a/a_{\rm eq}$可整理为
\begin{equation}\label{eq:astro-meszaros}
 D_{yy}+\frac{2+3y}{2y(1+y)}D_y-\frac{3}{2y(1+y)}D=0.
\end{equation}
一个解是$D_1=1+3y/2$，另一个可由降阶得到
\[
 D_2=\left(1+\frac32y\right)
 \ln\frac{\sqrt{1+y}+1}{\sqrt{1+y}-1}-3\sqrt{1+y}.
\]
在$y\ll1$时两解约为常数与$\ln(4/y)-3$；$y\gg1$时增长解趋于$a$。进入视界后的辐射势源与初始匹配通常造成冷暗物质的缓慢对数增长；原课的“停滞”是相对于物质时期$a$增长的概括，不是精确地处处常数。热或暖暗物质还会自由流动擦除一定尺度以下的起伏，不能把所有暗物质都当成无任何小尺度截止的CDM。

\originalsection{视界外的密度增长为何依赖规范}
在视界外，牛顿Poisson方程与瞬时引力解释不能直接外推。改变时间切片$t\mapsto t+T(\mathbf x)$会把$\delta\rho$变为$\delta\rho-\dot{\bar\rho}T$，所以$\delta$自身依赖规范。课程所写的辐射时期$\delta\propto a^2$、物质时期$\delta\propto a$可用于合适的共动密度变量；它们不表示远处区域在超光速交换信息，也不表示任意规范的密度都增长。

以$c=1$的共形Newton规范为例，$\dd s^2=a^2[-(1+2\Psi)\dd\eta^2+(1-2\Phi)\dd\mathbf x^2]$，$\mathcal H=a'/a$，$\theta=\ii\mathbf k\cdot\mathbf v$。无压强冷物质的守恒给
\[
 \delta_c'=-\theta_c+3\Phi',\qquad
 \theta_c'+\mathcal H\theta_c=k^2\Psi.
\]
因此辐射与势的时间变化不能通过把牛顿式的$4\pi G\rho$换成$32\pi G\rho/3$而全部处理。无剪切的自由辐射还含$\delta_\gamma'=-4\theta_\gamma/3+4\Phi'$、$\theta_\gamma'=k^2(\delta_\gamma/4+\Psi)$；有Thomson散射则加入拖曳，自由流中微子还需高阶矩。这里仅用于说明遗漏的项，未声称完成一般多组分Boltzmann数值演化。符号与规范限制另核对\href{https://arxiv.org/abs/astro-ph/9506072}{Ma与Bertschinger（1995）}的守恒式及CDM式，原论文不收录到源码包。

\originalsection{增长对星系形成的启示与限度}
只让重子从$z\sim1100$开始、以$D\propto a$增长，放大倍数只有约$10^3$。CMB的温度起伏约$10^{-5}$说明一个纯重子、标准初始条件的简单模型难以产生足够的晚期结构，已有冷暗物质势阱改善这一问题。但温度起伏不等于所有尺度上的物质$\delta$，还涉及传递函数、声学效应和投影。原课“没有CDM就不可能有结构”应限定在所讨论的引力、初始谱和背景模型中，不是关于任何替代理论的无条件定理。下一章把不同尺度的增长转成可与巡天比较的统计量。
